\chapter{Concentration of Measure}
\begin{center}\itshape
A function of many independent variables that does not depend too much on any one of them is nearly constant.\\
This simple principle, made quantitative, explains why high-dimensional spheres are thin, why random projections preserve distances,\\
and why the longest increasing subsequence of a random permutation hardly fluctuates.
\end{center}

In his lectures on functional analysis of 1922, Paul L\'evy proved that on a high-dimensional sphere, a set of half the area has almost the whole sphere within a short distance of it: the
measure ``concentrates'' around any equator. For half a century this looked like a curiosity of geometry. Then in 1971 Vitali Milman used it to give a new
proof of Dvoretzky's theorem on almost Euclidean sections of convex bodies, and saw that the same phenomenon is everywhere: in Gaussian space, on the cube,
in products of arbitrary probability spaces. Its master was Michel Talagrand, whose 1995 memoir on ``concentration of measure and isoperimetric inequalities
in product spaces'' gave inequalities strong enough to handle combinatorial problems that had resisted everything else. The Abel Prize of 2024 honoured him
``for his groundbreaking contributions to probability theory and functional analysis, with outstanding applications in mathematical physics and statistics''.

Tome~3 met the first instance: the bounded difference inequality, proved with martingales. This chapter sharpens it, adds the variance version of
Efron and Stein, proves the Gaussian concentration inequality through the Ornstein--Uhlenbeck semigroup and the logarithmic Sobolev inequality, and ends with
Talagrand's convex distance inequality. The applications: random projections that compress data without distorting it (Johnson and Lindenstrauss), and Ulam's
problem on increasing subsequences.

\begin{reminder}[from Tome 3, Tome 6 and Chapter 6]
Hoeffding's lemma: if $Y$ takes values in an interval of length $c$ and $\mathbb EY=0$, then $\mathbb Ee^{\lambda Y}\le e^{\lambda^2c^2/8}$ (Tome~3, Lemma~4.8, rescaled from $[0,1]$).
The bounded difference inequality (Tome~3, Theorem~10.25): if changing the $i$-th coordinate changes $f$ by at most $c_i$, then
$P(|F-\mathbb EF|\ge t)\le2e^{-t^2/2C}$ with $C=\sum c_i^2$. Fubini's theorem (Tome~6, Theorem~8.6) lets us integrate out some coordinates of a function of independent
variables. The Gaussian integration by parts $\mathbb E\,Zf(Z)=\mathbb Ef'(Z)$ is Lemma~\ref{lem:stein}.
\end{reminder}

Throughout, $X=(X_1,\dots,X_n)$ has independent coordinates, $X_i$ with law $\mu_i$ on a measurable space $E_i$, and $f:E_1\times\dots\times E_n\to\R$ is measurable. For
$x\in\prod E_i$ we write $x^{(i)}$ for $x$ with its $i$-th coordinate replaced: $X^{(i)}=(X_1,\dots,X_i',\dots,X_n)$, where $X'$ is an independent copy of $X$.

\section{Variance: the Efron--Stein inequality}

\begin{theorem}[Efron--Stein]\label{thm:efronstein}
If $\mathbb Ef(X)^2<\infty$, then
\[
\operatorname{Var}f(X)\le\frac12\sum_{i=1}^n\mathbb E\bigl[(f(X)-f(X^{(i)}))^2\bigr].
\]
\end{theorem}
\begin{proof}
For $0\le k\le n$ let
\[
g_k(x_1,\dots,x_k)=\int f(x_1,\dots,x_k,y_{k+1},\dots,y_n)\,\mu_{k+1}(\dd y_{k+1})\cdots\mu_n(\dd y_n),
\]
so that $g_n=f$ and $g_0=\mathbb Ef(X)$ (Fubini). Put $\Delta_k=g_k(X_1,\dots,X_k)-g_{k-1}(X_1,\dots,X_{k-1})$. For $j<k$, integrating $\Delta_k$ in $x_k$ with the earlier coordinates fixed gives $0$, so
$\mathbb E\Delta_j\Delta_k=0$, and $\operatorname{Var}f(X)=\mathbb E\bigl(\sum_k\Delta_k\bigr)^2=\sum_k\mathbb E\Delta_k^2$.

Now let $\bar f_k$ be the function $f$ with its $k$-th coordinate integrated out: $\bar f_k(x)=\int f(x_1,\dots,z,\dots,x_n)\,\mu_k(\dd z)$. Then
\[
\Delta_k=\int\bigl(f-\bar f_k\bigr)(X_1,\dots,X_k,y_{k+1},\dots,y_n)\,\prod_{j>k}\mu_j(\dd y_j),
\]
and by the Cauchy--Schwarz inequality,
$\mathbb E\Delta_k^2\le\mathbb E\bigl[(f-\bar f_k)(X)^2\bigr]$. Finally, for fixed values of the other coordinates, $\int(f-\bar f_k)^2\dd\mu_k$ is the variance of $f$ in the $k$-th
coordinate, which is
\[
\frac12\iint\bigl(f(\dots,z,\dots)-f(\dots,z',\dots)\bigr)^2\mu_k(\dd z)\,\mu_k(\dd z').
\]
Integrating over the other coordinates gives $\frac12\mathbb E(f(X)-f(X^{(k)}))^2$.
\end{proof}

For $f(x)=x_1+\dots+x_n$ this is an equality (Exercise~\ref{ex:esequality}). In general it says: the variance is at most the sum of the sensitivities to single coordinates.
Since $f(X)-f(X^{(i)})$ is symmetric in law, $\frac12\mathbb E(f-f^{(i)})^2=\mathbb E[(f-f^{(i)})_+^2]$, where $f^{(i)}=f(X^{(i)})$, and only the coordinates whose change can \emph{lower}
$f$ count. That makes the inequality powerful for maxima.

\begin{proposition}[increasing subsequences have small variance]\label{prop:lisvar}
Let $X_1,\dots,X_n$ be i.i.d.\ uniform on $[0,1]$ and $L_n$ the length of their longest increasing subsequence. Then $\sqrt n\le\mathbb EL_n\le2e\sqrt n+1$ and
$\operatorname{Var}L_n\le\mathbb EL_n$.
\end{proposition}
\begin{proof}
Ties have probability $0$, and the order of the $X_i$ is a uniform random permutation. \emph{Mean.} By the Erd\H os--Szekeres theorem (Tome~3, Theorem~A.15), a sequence of
$n$ distinct numbers with no increasing subsequence longer than $L$ and no decreasing one longer than $D$ has $n\le LD$. So $L_n+D_n\ge2\sqrt{L_nD_n}\ge2\sqrt n$; reversing
the sequence shows that $L_n$ and $D_n$ have the same law, so $\mathbb EL_n\ge\sqrt n$. The upper bound is Exercise~\ref{ex:lismean}.
\emph{Variance.} Fix a rule choosing one longest increasing subsequence $I$ of $X$. If $i\notin I$, replacing $X_i$ leaves $I$ increasing, so $L_n^{(i)}\ge L_n$. If $i\in I$, the
replacement lowers $L_n$ by at most $1$. So $\sum_i(L_n-L_n^{(i)})_+^2\le|I|=L_n$, and Theorem~\ref{thm:efronstein} gives $\operatorname{Var}L_n\le\mathbb EL_n$.
\end{proof}

So $L_n$, of order $\sqrt n$, fluctuates by at most $O(n^{1/4})$. This was the question asked by Stanis\l aw Ulam in 1961, and the story ends at the close of
the chapter.

\section{Exponential bounds: McDiarmid's inequality}

\begin{theorem}[McDiarmid]\label{thm:mcdiarmid}
Suppose $|f(x)-f(x^{(i)})|\le c_i$ for all $x$, all $i$ and all replacements of the $i$-th coordinate, and let $C=\sum_ic_i^2$. Then for $t>0$,
\[
P\bigl(f(X)-\mathbb Ef(X)\ge t\bigr)\le e^{-2t^2/C},\qquad P\bigl(|f(X)-\mathbb Ef(X)|\ge t\bigr)\le2e^{-2t^2/C}.
\]
\end{theorem}
\begin{proof}
With $g_k$ and $\Delta_k$ as in the previous proof ($f$ is bounded, since changing coordinates one at a time moves it by at most $\sum c_i$), fix $x_1,\dots,x_{k-1}$. As a
function of $z$, $g_k(x_1,\dots,x_{k-1},z)$ is an average over the later coordinates of functions whose values, as $z$ varies, lie in an interval of length $c_k$; so it
takes values in an interval of length $c_k$, and its $\mu_k$-mean is $g_{k-1}(x_1,\dots,x_{k-1})$. By Hoeffding's lemma,
\[
\int e^{\lambda(g_k(x_1,\dots,x_{k-1},z)-g_{k-1}(x_1,\dots,x_{k-1}))}\,\mu_k(\dd z)\le e^{\lambda^2c_k^2/8}.
\]
Integrating out $x_n$, then $x_{n-1}$, and so on (Fubini), $\mathbb Ee^{\lambda(f(X)-\mathbb Ef(X))}=\mathbb Ee^{\lambda(\Delta_1+\dots+\Delta_n)}\le e^{\lambda^2C/8}$. Markov's inequality with
$\lambda=4t/C$ gives $e^{-\lambda t+\lambda^2C/8}=e^{-2t^2/C}$. Apply this to $-f$ for the lower tail.
\end{proof}

Compared with Tome~3, the exponent improves by a factor of $4$, because Hoeffding's lemma sees that each increment lies in an interval of length $c_k$, not in
$[-c_k,c_k]$. For the $m=10^6$ balls in $n=10^6$ bins of Tome~3, the probability that the number of empty bins deviates from its mean by $5000$ drops from at most
$2e^{-12.5}$ to at most $2e^{-50}$.

\section{Gaussian concentration}

Let $Z$ be a standard Gaussian vector in $\R^n$ and $\gamma$ its law. We work with functions in the class $\mathcal C$ of $C^2$ functions on $\R^n$ whose values and first and
second partial derivatives are bounded by $Ce^{c|x|}$ for some constants: Gaussian expectations of such functions are finite, and we may differentiate under them
(dominated convergence). The \emph{Ornstein--Uhlenbeck semigroup} is
\[
P_th(x)=\mathbb E\,h\bigl(e^{-t}x+\sqrt{1-e^{-2t}}\,Z\bigr),\qquad t\ge0.
\]
It interpolates between $P_0h=h$ and $P_\infty h=\mathbb Eh(Z)$: for each $x$, $P_th(x)\to\mathbb Eh(Z)$ as $t\to\infty$ by dominated convergence.

\begin{lemma}[Ornstein--Uhlenbeck calculus]\label{lem:ou}
For $h,g\in\mathcal C$ and $t\ge0$, with $Lh=\Delta h-x\cdot\nabla h$:
\begin{enumerate}[label=(\roman*)]
\item $\mathbb EP_th(Z)=\mathbb Eh(Z)$;
\item $\nabla P_th=e^{-t}P_t(\nabla h)$ (coordinatewise);
\item $\frac{\partial}{\partial t}P_th=LP_th$;
\item $\mathbb E[g(Z)\,Lh(Z)]=-\mathbb E[\nabla g(Z)\cdot\nabla h(Z)]$.
\end{enumerate}
\end{lemma}
\begin{proof}
Write $\sigma_t=\sqrt{1-e^{-2t}}$ and $Y=e^{-t}x+\sigma_tZ$. (i) If $Z'$ is an independent copy of $Z$, then $e^{-t}Z'+\sigma_tZ$ is centred Gaussian with covariance $(e^{-2t}+\sigma_t^2)I=I$.
(ii) Differentiate under the expectation. (iv) Apply Lemma~\ref{lem:stein}(i) in the $i$-th coordinate to $u=g\,\partial_ih$ (with the others fixed, then Fubini):
$\mathbb E[\partial_ig\,\partial_ih+g\,\partial_i^2h]=\mathbb E[Z_ig\,\partial_ih]$; sum over $i$. (iii) By the chain rule, $\frac{\partial}{\partial t}P_th(x)=\mathbb E\bigl[\nabla h(Y)\cdot(-e^{-t}x+\sigma_t'Z)\bigr]$.
Gaussian integration by parts in $Z$ turns $\mathbb E[\partial_ih(Y)Z_i]$ into $\sigma_t\mathbb E[\partial_i^2h(Y)]$, and $\sigma_t\sigma_t'=e^{-2t}$. So
$\frac{\partial}{\partial t}P_th=e^{-2t}P_t(\Delta h)-e^{-t}x\cdot P_t(\nabla h)$, which by (ii) is $\Delta P_th-x\cdot\nabla P_th$.
\end{proof}

The semigroup trades time for smoothness: it starts from $h$, ends at a constant, and along the way (ii) shrinks gradients at rate $e^{-t}$. Integrating the change
over time gives two inequalities.

\begin{theorem}[Gaussian Poincar\'e and logarithmic Sobolev inequalities]\label{thm:gausslsi}
For $f\in\mathcal C$,
\[
\operatorname{Var}f(Z)\le\mathbb E|\nabla f(Z)|^2.
\]
For $h=e^G$ with $G\in C^2$ having bounded first and second derivatives, writing $\operatorname{Ent}(h)=\mathbb E[h\log h]-\mathbb Eh\log\mathbb Eh$ (all at $Z$),
\[
\operatorname{Ent}(h)\le\frac12\,\mathbb E\Bigl[\frac{|\nabla h|^2}h\Bigr].
\]
\end{theorem}
\begin{proof}
Let $\phi$ be $C^2$ on the range of the functions involved, and $u_t=P_th$. By Lemma~\ref{lem:ou}(iii) and (iv),
\[
\frac{\dd}{\dd t}\mathbb E\phi(u_t(Z))=\mathbb E\bigl[\phi'(u_t)Lu_t\bigr]=-\mathbb E\bigl[\phi''(u_t)|\nabla u_t|^2\bigr].
\]
Integrating from $0$ to $\infty$ and using $u_t\to\mathbb Eh(Z)$,
\begin{equation}\label{eq:ouint}
\mathbb E\phi(h(Z))-\phi(\mathbb Eh(Z))=\int_0^\infty\mathbb E\bigl[\phi''(u_t)|\nabla u_t|^2\bigr]\dd t.
\end{equation}
\emph{Poincar\'e.} Take $h=f$ and $\phi(u)=u^2$. By (ii) and Jensen, $|\nabla u_t|^2=e^{-2t}|P_t\nabla f|^2\le e^{-2t}P_t(|\nabla f|^2)$, whose expectation is $e^{-2t}\mathbb E|\nabla f(Z)|^2$
by (i). So the right side of~\eqref{eq:ouint} is at most $\int_0^\infty2e^{-2t}\dd t\,\mathbb E|\nabla f|^2$.

\emph{Log-Sobolev.} Here $h>0$, $h\in\mathcal C$, and we take $\phi(u)=u\log u$, $\phi''(u)=\frac1u$. By (ii) and the Cauchy--Schwarz inequality for the expectation defining $P_t$,
\[
|\nabla u_t|=e^{-t}|P_t\nabla h|\le e^{-t}P_t|\nabla h|=e^{-t}P_t\Bigl(\frac{|\nabla h|}{\sqrt h}\sqrt h\Bigr)\le e^{-t}\Bigl(P_t\frac{|\nabla h|^2}h\Bigr)^{1/2}(P_th)^{1/2}.
\]
So $\frac{|\nabla u_t|^2}{u_t}\le e^{-2t}P_t\frac{|\nabla h|^2}h$, and~\eqref{eq:ouint} with (i) gives $\operatorname{Ent}(h)\le\int_0^\infty e^{-2t}\dd t\,\mathbb E\frac{|\nabla h|^2}h$.
\end{proof}

The logarithmic Sobolev inequality is due to Leonard Gross (1975), and the semigroup proof is in the spirit of Bakry and \'Emery. Its strength is that it
tensorizes without loss and turns into concentration by a two-line argument found by Ira Herbst in an unpublished letter.

\begin{theorem}[Gaussian concentration]\label{thm:gaussconc}
Let $F:\R^n\to\R$ be $L$-Lipschitz for the Euclidean norm. Then for $t>0$,
\[
P\bigl(F(Z)-\mathbb EF(Z)\ge t\bigr)\le e^{-t^2/2L^2},\qquad\operatorname{Var}F(Z)\le L^2.
\]
\end{theorem}
\begin{proof}
By scaling, $L=1$. \emph{Smooth case.} Suppose first that $F$ is $C^2$ with bounded second derivatives; being $1$-Lipschitz, $|\nabla F|\le1$. The Poincar\'e inequality gives
$\operatorname{Var}F\le1$. For $\lambda>0$ let $H(\lambda)=\mathbb Ee^{\lambda F(Z)}$. The log-Sobolev inequality for $h=e^{\lambda F}$, for which $|\nabla h|^2/h=\lambda^2|\nabla F|^2h\le\lambda^2h$, reads
\[
\lambda H'(\lambda)-H(\lambda)\log H(\lambda)\le\frac{\lambda^2}2H(\lambda)
\]
(Herbst's argument). Divide by $\lambda^2H$: the function $K(\lambda)=\frac1\lambda\log H(\lambda)$ satisfies $K'(\lambda)\le\frac12$, and $K(\lambda)\to\frac{H'(0)}{H(0)}=\mathbb EF$ as $\lambda\to0$. So
$\log\mathbb Ee^{\lambda(F-\mathbb EF)}\le\frac{\lambda^2}2$, and Markov's inequality with $\lambda=t$ gives $e^{-t^2/2}$.

\emph{General case.} Let $W$ be a standard Gaussian vector independent of $Z$ and $F_\varepsilon(x)=\mathbb EF(x+\varepsilon W)$. It is $1$-Lipschitz, $|F_\varepsilon-F|\le\varepsilon\mathbb E|W|\le\varepsilon\sqrt n$, and it
is $C^2$ with bounded second derivatives: differentiating the Gaussian density of $x+\varepsilon W$ twice,
$\partial_i\partial_jF_\varepsilon(x)=\varepsilon^{-2}\mathbb E[(F(x+\varepsilon W)-F(x))(W_iW_j-\delta_{ij})]$, bounded by $\varepsilon^{-1}\mathbb E[|W|\,|W_iW_j-\delta_{ij}|]$. So
$P(F-\mathbb EF\ge t)\le P(F_\varepsilon-\mathbb EF_\varepsilon\ge t-2\varepsilon\sqrt n)\le e^{-(t-2\varepsilon\sqrt n)^2/2}$ and $\operatorname{Var}F\le(1+\varepsilon\sqrt n)^2$ (Minkowski's inequality for the standard deviations of $F_\varepsilon$ and $F-F_\varepsilon$); let $\varepsilon\to0$.
\end{proof}

The bound does not depend on the dimension. Two consequences show what that means.

\begin{example}[the thin shell]\label{ex:thinshell}
$F(z)=|z|$ is $1$-Lipschitz, so $|Z|$ deviates from its mean by more than $t$ with probability at most $2e^{-t^2/2}$, while $\mathbb E|Z|^2=n$ and $\sqrt n-1\le\mathbb E|Z|\le\sqrt n$
(Exercise~\ref{ex:thinshellmean}). A Gaussian vector in dimension $10^6$ has a length within $3$ of its mean, which is between $999$ and $1000$, with probability at least $0.97$: its
law is a thin shell, not a fuzzy ball (Figure~\ref{fig:thinshell}).
\end{example}

\begin{figure}[ht]
\centering\fig{ch08-thinshell}
\caption{Histograms of $|Z|-\sqrt n$ for a standard Gaussian vector $Z$ in dimension $n=2$, $20$, $200$ and $2000$ ($10^5$ samples each). The mean tends to $0$ (it is about
$-\frac1{4\sqrt n}$, see Exercise~\ref{ex:thinshellmean}) while the standard deviation stays near $\frac1{\sqrt2}\approx0.71$, as Theorem~\ref{thm:gaussconc} predicts: the
length of $Z$ is $\sqrt n$ up to an error of order one, whatever the dimension.}
\label{fig:thinshell}
\end{figure}

\begin{example}[maxima of Gaussian vectors]\label{ex:gaussmax}
Let $G=AZ$ be a centred Gaussian vector in $\R^m$ whose coordinates have variances $\sum_jA_{ij}^2\le1$. Then $F(z)=\max_i(Az)_i$ is $1$-Lipschitz (each $(Az)_i$ is, by
Cauchy--Schwarz), so $\max_iG_i$ fluctuates around its mean by $O(1)$, \emph{whatever the correlations}, while its mean may grow like $\sqrt{2\log m}$
(Exercise~\ref{ex:maxmean}).
\end{example}

\section{Random projections: Johnson--Lindenstrauss}

A data set of $N$ points in $\R^d$, with $d$ huge (documents as word counts, images as pixels, users as ratings), must be stored and searched by distance. William Johnson
and Joram Lindenstrauss proved in 1984, as a lemma in a paper about extending Lipschitz maps, that a \emph{random} linear map to dimension $k=O(\varepsilon^{-2}\log N)$
preserves all pairwise distances up to a factor $1\pm\varepsilon$. The dimension $d$ does not appear.

\begin{lemma}[chi-square tails]\label{lem:chisq}
Let $Q=Z_1^2+\dots+Z_k^2$ with $Z_i$ independent standard Gaussians, and $0<\varepsilon<1$. Then
\[
P\bigl(Q\ge(1+\varepsilon)k\bigr)\le e^{-\frac k2(\frac{\varepsilon^2}2-\frac{\varepsilon^3}3)},\qquad P\bigl(Q\le(1-\varepsilon)k\bigr)\le e^{-\frac k2\cdot\frac{\varepsilon^2}2}.
\]
\end{lemma}
\begin{proof}
$\mathbb Ee^{\theta Z^2}=(1-2\theta)^{-1/2}$ for $\theta<\frac12$ (a Gaussian integral). Upper tail: with $\theta=\frac\varepsilon{2(1+\varepsilon)}$, Markov gives
$P(Q\ge(1+\varepsilon)k)\le(1-2\theta)^{-k/2}e^{-\theta(1+\varepsilon)k}=\bigl((1+\varepsilon)e^{-\varepsilon}\bigr)^{k/2}$, and $\log(1+\varepsilon)-\varepsilon\le-\frac{\varepsilon^2}2+\frac{\varepsilon^3}3$. Lower tail: with
$\theta=\frac\varepsilon{2(1-\varepsilon)}$, $P(Q\le(1-\varepsilon)k)\le\mathbb Ee^{-\theta Q}e^{\theta(1-\varepsilon)k}=\bigl((1-\varepsilon)e^{\varepsilon}\bigr)^{k/2}$, and $\log(1-\varepsilon)+\varepsilon\le-\frac{\varepsilon^2}2$.
\end{proof}

\begin{theorem}[Johnson--Lindenstrauss]\label{thm:jl}
Let $x_1,\dots,x_N\in\R^d$, $0<\varepsilon<1$, and $k>\frac{4\log N}{\varepsilon^2/2-\varepsilon^3/3}$. Let $A$ be a $k\times d$ matrix of independent $N(0,\frac1k)$ entries. Then with probability
at least $1-N^2e^{-\frac k2(\frac{\varepsilon^2}2-\frac{\varepsilon^3}3)}>0$, for all $i,j$,
\[
(1-\varepsilon)|x_i-x_j|^2\le|Ax_i-Ax_j|^2\le(1+\varepsilon)|x_i-x_j|^2.
\]
\end{theorem}
\begin{proof}
For a fixed unit vector $u$, the coordinates of $\sqrt kAu$ are independent standard Gaussians (each is $\sum_jg_{ij}u_j$ with $g_{ij}$ i.i.d.\ $N(0,1)$ and $\sum u_j^2=1$),
so $k|Au|^2$ has the law of $Q$ in Lemma~\ref{lem:chisq}, and $|Au|^2\notin[1-\varepsilon,1+\varepsilon]$ with probability at most $2e^{-\frac k2(\frac{\varepsilon^2}2-\frac{\varepsilon^3}3)}$. Apply this to
$u=\frac{x_i-x_j}{|x_i-x_j|}$ for the $\binom N2<\frac{N^2}2$ pairs and use the union bound. The probability is positive by the choice of $k$.
\end{proof}

For a million points and $\varepsilon=0.1$ the theorem asks for $k\ge11\,842$, independent of $d$: a database of a million documents described by ten million word counts
can be compressed a thousandfold and still be searched by distance. The constant is not the point; the logarithm is. In practice one uses $\pm1$ entries or sparse
matrices, which are faster and satisfy the same lemma (Achlioptas, 2003), and numerical evaluation shows how much smaller $k$ can be in practice.

\section{Talagrand's convex distance}

Bounded differences treat all coordinates alike. Many functions, like the length of a longest increasing subsequence, change by one when \emph{some} coordinate changes,
but only a few coordinates matter for any given input. Talagrand found the right distance for this.

For a product space $\Omega=\prod_iE_i$ (in this section each $E_i$ is countable, so that everything is measurable), $x\in\Omega$ and $A\subseteq\Omega$, the \emph{convex distance} is
\[
d_T(x,A)=\sup_{\alpha}\ \inf_{y\in A}\ \sum_{i:\,x_i\ne y_i}\alpha_i,
\]
the supremum over $\alpha\in\R^n$ with $\alpha_i\ge0$ and $\sum\alpha_i^2=1$. With $\alpha=\frac1{\sqrt n}(1,\dots,1)$ it is at least $\frac1{\sqrt n}$ times the Hamming distance; the freedom to
choose $\alpha$ \emph{depending on $x$} is the point.

\begin{lemma}[convex distance as a Euclidean distance]\label{lem:dtgeom}
Let $U_A(x)$ be the set of $s\in\{0,1\}^n$ such that some $y\in A$ has $y_i=x_i$ whenever $s_i=0$, and $V_A(x)$ its convex hull. If $A\ne\emptyset$, then
$d_T(x,A)=\min\{|v|:v\in V_A(x)\}$.
\end{lemma}
\begin{proof}
For a given $\alpha$, $\inf_{y\in A}\sum_{x_i\ne y_i}\alpha_i=\min_{s\in U_A(x)}\alpha\cdot s=\min_{v\in V_A(x)}\alpha\cdot v$ (the indicator of $\{x_i\ne y_i\}$ belongs to $U_A(x)$, every $s\in U_A(x)$ dominates one
such indicator, and a linear function is minimized on a convex hull at a vertex). If $v\in V_A(x)$, $\alpha\cdot v\le|v|$, so $d_T\le\min|v|$. Conversely let $v^*$ be the point of the compact
convex set $V_A(x)$ closest to $0$. If $v^*=0$ both sides vanish. Otherwise, for $v\in V_A(x)$ and $0<s\le1$, $|v^*+s(v-v^*)|^2\ge|v^*|^2$ gives $v\cdot v^*\ge|v^*|^2$;
so $\alpha=v^*/|v^*|$, which has non-negative coordinates, gives $\min_{v\in V}\alpha\cdot v\ge|v^*|$.
\end{proof}

\begin{theorem}[Talagrand's convex distance inequality]\label{thm:talagrand}
For every $A\subseteq\Omega$ with $P(X\in A)>0$,
\[
\mathbb E\,e^{d_T(X,A)^2/4}\le\frac1{P(X\in A)}.\qquad\text{Consequently}\qquad P(X\in A)\,P\bigl(d_T(X,A)\ge t\bigr)\le e^{-t^2/4}.
\]
\end{theorem}
\begin{proof}
The consequence follows by Markov's inequality. We prove the first inequality by induction on $n$, writing $P$ for the law of $X$ on each product space.

\emph{An inequality of numbers.} For $0\le r\le1$,
\[
\min_{0\le\lambda\le1}r^{-\lambda}e^{(1-\lambda)^2/4}\le2-r.
\] If $r<e^{-1/2}$, take $\lambda=0$: $e^{1/4}<1.29<2-e^{-1/2}$. Otherwise take
$\lambda=1+2\log r\in[0,1]$; the value is $\exp(-u-u^2)$ with $u=\log r\in[-\frac12,0]$, and $\psi(u)=\log(2-e^u)+u+u^2$ satisfies $\psi(0)=0$, $\psi'(0)=0$ and
$\psi''(u)=2-\frac{2e^u}{(2-e^u)^2}\ge0$ for $u\le0$, so $\psi'\le0$ and $\psi\ge0$ on $[-\frac12,0]$: $e^{-u-u^2}\le2-e^u$.

\emph{$n=1$.} Here $d_T(x,A)$ is $0$ if $x\in A$ and $1$ otherwise. With $p=P(A)$: $p+(1-p)e^{1/4}\le\frac1p$, since $\frac1p-p=(1-p)\frac{1+p}p\ge2(1-p)$.

\emph{From $n$ to $n+1$.} Write points of $\Omega\times E$ as $(x,\omega)$. Let $B=\{x:(x,\omega)\in A\text{ for some }\omega\}$ and $A_\omega=\{x:(x,\omega)\in A\}$. If $s\in U_{A_\omega}(x)$ then $(s,0)\in U_A(x,\omega)$
(the witness $y$ in $A_\omega$ gives the witness $(y,\omega)$ in $A$); if $s\in U_B(x)$ then $(s,1)\in U_A(x,\omega)$. Hence for $v\in V_{A_\omega}(x)$, $w\in V_B(x)$ and $\lambda\in[0,1]$, the point
$(\lambda v+(1-\lambda)w,\,1-\lambda)$ lies in $V_A(x,\omega)$, and by Lemma~\ref{lem:dtgeom} and the convexity of $|\cdot|^2$,
\[
d_T((x,\omega),A)^2\le(1-\lambda)^2+\lambda\,d_T(x,A_\omega)^2+(1-\lambda)\,d_T(x,B)^2.
\]
Fix $\omega$ with $P(A_\omega)>0$. Exponentiate, integrate in $x$, use H\"older's inequality with exponents $\frac1\lambda$, $\frac1{1-\lambda}$, and the induction hypothesis for $A_\omega$ and $B$:
\[
\int e^{d_T((x,\omega),A)^2/4}P(\dd x)\le e^{(1-\lambda)^2/4}P(A_\omega)^{-\lambda}P(B)^{-(1-\lambda)}=\frac1{P(B)}\,r^{-\lambda}e^{(1-\lambda)^2/4},
\]
with $r=P(A_\omega)/P(B)\le1$. Minimize over $\lambda$: the integral is at most $\frac{2-r}{P(B)}$, a bound that also holds when $P(A_\omega)=0$ (take $\lambda=0$). Integrate over $\omega$; since
$\int P(A_\omega)\,P(\dd\omega)=P(A)$, with $a=P(A)/P(B)\in(0,1]$,
\[
\mathbb Ee^{d_T(X,A)^2/4}\le\frac{2-a}{P(B)}=\frac{a(2-a)}{P(A)}\le\frac1{P(A)} ,
\]
since $a(2-a)=1-(1-a)^2\le1$.
\end{proof}

How is it used? The function only needs short \emph{certificates}: if $f(x)\ge s$ can be witnessed by a few coordinates, the convex distance controls how far below $s$
a nearby point can go.

\begin{theorem}[concentration of increasing subsequences]\label{thm:lisconc}
Let $L_n$ be the length of the longest increasing subsequence of a uniform random permutation of $\{1,\dots,n\}$, and $M_n$ a median of $L_n$. For $t>0$,
\[
P(L_n\ge M_n+t)\le2e^{-t^2/4(M_n+t)},\qquad P(L_n\le M_n-t)\le2e^{-t^2/4M_n}.
\]
\end{theorem}
\begin{proof}
\emph{The certificate.} Let $X_1,\dots,X_n$ be independent with values in a countable totally ordered set, $\ell(x)$ the length of the longest strictly increasing subsequence of
$x$, and $a\ge0$. Let $A=\{y:\ell(y)\le a\}$. If $\ell(x)=\ell$, let $I$ be the index set of a longest increasing subsequence of $x$ and $\alpha=\mathbf 1_I/\sqrt\ell$. For $y\in A$, the indices of $I$
where $y$ agrees with $x$ carry an increasing subsequence of $y$, so there are at most $a$ of them, and $\sum_{x_i\ne y_i}\alpha_i\ge\frac{\ell-a}{\sqrt\ell}$. As $\frac{\ell-a}{\sqrt\ell}$ increases with $\ell$,
$\ell(x)\ge a+t$ implies $d_T(x,A)\ge\frac t{\sqrt{a+t}}$, and Theorem~\ref{thm:talagrand} gives
\[
P(\ell(X)\le a)\,P(\ell(X)\ge a+t)\le e^{-t^2/4(a+t)}.
\]
\emph{Back to permutations.} Take $X_i$ uniform on $\{1,\dots,m\}$. With probability at least $1-\binom n2\frac1m$ the values are distinct, and then their relative order is a uniform permutation;
so as $m\to\infty$ the law of $\ell(X)$ tends to the law of $L_n$, and the inequality holds for $L_n$. With $a=M_n$, $P(L_n\le M_n)\ge\frac12$ gives the first bound; with
$a=M_n-t$, $P(L_n\ge M_n)\ge\frac12$ gives the second.
\end{proof}

Ulam asked in 1961 how $L_n$ behaves. John Hammersley proved in 1972 that $L_n/\sqrt n$ converges in probability to a constant $c$, and in 1977 Benjamin Logan and Larry Shepp, and
independently Anatoly Vershik and Sergei Kerov, showed that $c=2$. With $M_n\approx2\sqrt n$, Theorem~\ref{thm:lisconc} says that the fluctuations are at most of order
$n^{1/4}$, which is what Talagrand obtained in 1995. The truth is smaller still: in 1999 Jinho Baik, Percy Deift and Kurt Johansson proved that
$(L_n-2\sqrt n)/n^{1/6}$ converges to the Tracy--Widom law of the largest eigenvalue of a random Hermitian matrix, a connection nobody had expected (Figure~\ref{fig:lis}).

\begin{figure}[ht]
\centering\fig{ch08-lis}
\caption{A uniform random permutation of $\{1,\dots,300\}$ drawn as the points $(i,\sigma(i))$, with a longest increasing subsequence in red; its length is $29$,
compared with $2\sqrt{300}\approx34.6$ (the correction of order $n^{1/6}$ is negative). Concentration says that every random permutation of this size has an
increasing path of nearly the same length.}
\label{fig:lis}
\end{figure}

\begin{think}
In the proof of Theorem~\ref{thm:lisconc}, the certificate of $\ell(x)\ge s$ has size $s$, and the bound depends on $s$, not on $n$. What would the bounded difference inequality
give for $L_n$? For which other functions (the size of a largest clique in a random graph, the length of a shortest travelling salesman tour through random points) can
you find small certificates?
\end{think}

\section{Exercises}

\begin{exercise}[$\star$]\label{ex:esequality}
Show that the Efron--Stein inequality is an equality for $f(x)=x_1+\dots+x_n$ (with square-integrable coordinates), and that it can be strict for $f(x)=x_1x_2$.
\end{exercise}

\begin{exercise}[$\star$]\label{ex:lismean}
Show that $P(L_n\ge k)\le\binom nk\frac1{k!}\le\bigl(\frac{e\sqrt n}k\bigr)^{2k}$, and deduce $\mathbb EL_n\le2e\sqrt n+1$.
\end{exercise}

\begin{exercise}[$\star$]\label{ex:thinshellmean}
Show that $\sqrt n-1\le\mathbb E|Z|\le\sqrt n$ for a standard Gaussian vector $Z$ in $\R^n$.
\end{exercise}

\begin{exercise}[$\star\star$]\label{ex:maxmean}
Let $G_1,\dots,G_m$ be centred Gaussians with variances at most $1$ (not necessarily independent). Show that $\mathbb E\max_iG_i\le\sqrt{2\log m}$, and that equality holds asymptotically
for independent standard Gaussians: $\mathbb E\max_iZ_i\ge(1-\varepsilon)\sqrt{2\log m}$ for large $m$.
\end{exercise}

\begin{exercise}[$\star\star$]\label{ex:rademachersum}
Let $v_1,\dots,v_n$ be vectors in a normed space and $\varepsilon_i$ independent fair signs. Show that $\|\sum_i\varepsilon_iv_i\|$ deviates from its mean by $t$ with probability at most
$2e^{-t^2/2\sum\|v_i\|^2}$.
\end{exercise}

\begin{exercise}[$\star\star\star$]\label{ex:convexlip}
Let $X_1,\dots,X_n$ be independent, each taking finitely many values in $[0,1]$, and $F:[0,1]^n\to\R$ convex and $1$-Lipschitz for the Euclidean norm, with median $M$. Show that
$P(|F(X)-M|\ge t)\le4e^{-t^2/4}$. (First show that if $A\subseteq[0,1]^n$ is convex, then the Euclidean distance from $x$ to $A$ is at most $d_T(x,A)$.)
\end{exercise}

\section{Hints}

\begin{hint}{ex:esequality}
Both sides equal $\sum\operatorname{Var}X_i$. For $x_1x_2$ with independent fair signs, compute both sides.
\end{hint}
\begin{hint}{ex:lismean}
Each $k$-subset of positions carries an increasing subsequence with probability $\frac1{k!}$. Use $\binom nk\le\frac{n^k}{k!}$ and $k!\ge(k/e)^k$; for $k\ge2e\sqrt n$ the bound is at most $4^{-k}$.
\end{hint}
\begin{hint}{ex:thinshellmean}
$\mathbb E|Z|\le\sqrt{\mathbb E|Z|^2}$; and $(\mathbb E|Z|)^2=n-\operatorname{Var}|Z|\ge n-1$.
\end{hint}
\begin{hint}{ex:maxmean}
$e^{\lambda\mathbb E\max G_i}\le\mathbb Ee^{\lambda\max G_i}\le\sum_i\mathbb Ee^{\lambda G_i}$. For the lower bound, $P(\max Z_i<x)=\Phi(x)^m$; use~\eqref{eq:gausstail} at $x=(1-\varepsilon)\sqrt{2\log m}$.
\end{hint}
\begin{hint}{ex:rademachersum}
Changing $\varepsilon_i$ moves the norm by at most $2\|v_i\|$. Use Theorem~\ref{thm:mcdiarmid}.
\end{hint}
\begin{hint}{ex:convexlip}
Write the minimal $v\in V_A(x)$ as a convex combination $\sum_j\lambda_js_j$ of points of $U_A(x)$ with witnesses $y_j\in A$, and look at $y=\sum_j\lambda_jy_j$. Then apply Theorem~\ref{thm:talagrand} to
$A=\{F\le M\}$ and to $A=\{F\le M-t\}$.
\end{hint}

\section{Solutions}

\begin{solution}{ex:esequality}
\idea For a sum, each coordinate contributes its own variance, independently: there is nothing to lose.
\writeup For $f=\sum x_i$, $f(X)-f(X^{(i)})=X_i-X_i'$, so the right side is $\frac12\sum\mathbb E(X_i-X_i')^2=\sum\operatorname{Var}X_i=\operatorname{Var}f(X)$. For $f=x_1x_2$ with fair signs, $\operatorname{Var}(X_1X_2)=1$,
while $f-f^{(1)}=(X_1-X_1')X_2$ has $\frac12\mathbb E(\cdot)^2=\frac12\cdot2=1$ and likewise for $i=2$: the right side is $2$.
\end{solution}

\begin{solution}{ex:lismean}
\idea An increasing subsequence of length $k$ lives on one of $\binom nk$ position sets, each of which is increasing with probability $\frac1{k!}$.
\writeup By the union bound over $k$-subsets of positions, each sorted with probability $\frac1{k!}$ (all orders of $k$ distinct i.i.d.\ values are equally likely),
$P(L_n\ge k)\le\binom nk\frac1{k!}\le\frac{n^k}{(k!)^2}\le\frac{n^k}{(k/e)^{2k}}=\bigl(\frac{e\sqrt n}k\bigr)^{2k}$, using $k!\ge(k/e)^k$ (from $e^k=\sum\frac{k^j}{j!}\ge\frac{k^k}{k!}$). For $k\ge2e\sqrt n$ this is at most $4^{-k}$. So
$\mathbb EL_n=\sum_{k\ge1}P(L_n\ge k)\le2e\sqrt n+\sum_{k\ge1}4^{-k}\le2e\sqrt n+1$.
\end{solution}

\begin{solution}{ex:thinshellmean}
\idea The mean of the length is squeezed between the root mean square, by Jensen, and the root mean square minus the variance, which concentration bounds by $1$.
\writeup $(\mathbb E|Z|)^2\le\mathbb E|Z|^2=n$. And $\operatorname{Var}|Z|\le1$ by Theorem~\ref{thm:gaussconc}, so $(\mathbb E|Z|)^2=n-\operatorname{Var}|Z|\ge n-1$ and $\mathbb E|Z|\ge\sqrt{n-1}\ge\sqrt n-1$. (In
fact $\mathbb E|Z|=\sqrt n-\frac1{4\sqrt n}+O(n^{-3/2})$.)
\end{solution}

\begin{solution}{ex:maxmean}
\idea The exponential moment of a maximum is at most the sum of the exponential moments: a union bound with a free parameter.
\writeup For $\lambda>0$, by Jensen, $e^{\lambda\mathbb E\max G_i}\le\mathbb Ee^{\lambda\max G_i}\le\sum_i\mathbb Ee^{\lambda G_i}\le me^{\lambda^2/2}$, so $\mathbb E\max G_i\le\frac{\log m}\lambda+\frac\lambda2$; take $\lambda=\sqrt{2\log m}$.
For the lower bound with independent $Z_i$, let $x=(1-\varepsilon)\sqrt{2\log m}$. By~\eqref{eq:gausstail}, $1-\Phi(x)\ge\frac x{1+x^2}\frac{e^{-x^2/2}}{\sqrt{2\pi}}=m^{-(1-\varepsilon)^2+o(1)}$, so
$P(\max Z_i<x)=\Phi(x)^m\le\exp(-m(1-\Phi(x)))\to0$. Since $\max Z_i\ge Z_1$, $\mathbb E\max Z_i\ge xP(\max Z_i\ge x)-\mathbb E|Z_1|$, which is $\ge(1-2\varepsilon)\sqrt{2\log m}$ for large $m$.
\end{solution}

\begin{solution}{ex:rademachersum}
\idea A norm of a sum is a function of independent signs with bounded differences.
\writeup Let $f(\varepsilon)=\|\sum\varepsilon_iv_i\|$. Flipping $\varepsilon_i$ changes the sum by $\pm2v_i$, so by the triangle inequality $f$ changes by at most $c_i=2\|v_i\|$. Theorem~\ref{thm:mcdiarmid}
gives $2e^{-2t^2/\sum4\|v_i\|^2}=2e^{-t^2/2\sum\|v_i\|^2}$.
\end{solution}

\begin{solution}{ex:convexlip}
\idea Convexity turns the combinatorial distance into a geometric one: averaging the witnesses produces a point of $A$ that is close coordinate by coordinate.
\writeup Let $A$ be convex and $v=\sum_j\lambda_js_j$ the point of $V_A(x)$ of minimal norm, with $s_j\in U_A(x)$ witnessed by $y_j\in A$, so $y_{j,i}=x_i$ whenever $s_{j,i}=0$. Then
$y=\sum_j\lambda_jy_j\in A$ and, since all coordinates lie in $[0,1]$, $|x_i-y_i|\le\sum_j\lambda_j|x_i-y_{j,i}|\le\sum_j\lambda_js_{j,i}=v_i$. So $\operatorname{dist}(x,A)\le|x-y|\le|v|=d_T(x,A)$.
Now let $A=\{F\le M\}$, convex. If $F(x)\ge M+t$, then every $y\in A$ has $|x-y|\ge F(x)-F(y)\ge t$, so $d_T(x,A)\ge t$, and Theorem~\ref{thm:talagrand} (the values form a countable
product space) gives $P(F\ge M+t)\le\frac{e^{-t^2/4}}{P(A)}\le2e^{-t^2/4}$. With $A=\{F\le M-t\}$, every $x$ with $F(x)\ge M$ has $d_T(x,A)\ge t$, so $P(F\le M-t)\cdot\frac12\le e^{-t^2/4}$. Add the two bounds.
\end{solution}

\section*{Sources and further reading}

\begin{itemize}
\item \emph{History.} P.~L\'evy, \emph{Le\c cons d'analyse fonctionnelle}, Gauthier-Villars, 1922 (the isoperimetric inequality on the sphere). V.~D.~Milman, ``A new proof of
A.~Dvoretzky's theorem on cross-sections of convex bodies'', \emph{Funkcional.\ Anal.\ i Prilo\v zen.}\ 5 (1971), 28--37. M.~Talagrand, ``Concentration of measure and isoperimetric
inequalities in product spaces'', \emph{Publ.\ Math.\ IH\'ES} 81 (1995), 73--205. The Abel Prize citation, 2024. B.~Efron and C.~Stein, ``The jackknife estimate of variance'',
\emph{Ann.\ Statist.}\ 9 (1981), 586--596 (the same Charles Stein as in Chapter~6). C.~McDiarmid, ``On the method of bounded differences'', in \emph{Surveys in Combinatorics 1989},
Cambridge University Press, 1989, 148--188. L.~Gross, ``Logarithmic Sobolev inequalities'', \emph{Amer.\ J.\ Math.}\ 97 (1975), 1061--1083. W.~B.~Johnson and J.~Lindenstrauss,
``Extensions of Lipschitz mappings into a Hilbert space'', \emph{Contemp.\ Math.}\ 26 (1984), 189--206. S.~Dasgupta and A.~Gupta, ``An elementary proof of a theorem of Johnson
and Lindenstrauss'', \emph{Random Structures Algorithms} 22 (2003), 60--65 (the constant used here). D.~Achlioptas, ``Database-friendly random projections'', \emph{J.\ Comput.\ System Sci.}\
66 (2003), 671--687.
\item \emph{Increasing subsequences.} J.~M.~Hammersley, ``A few seedlings of research'', \emph{Proc.\ Sixth Berkeley Symp.\ Math.\ Statist.\ Prob.}\ 1 (1972), 345--394. B.~F.~Logan and L.~A.~Shepp,
``A variational problem for random Young tableaux'', \emph{Adv.\ Math.}\ 26 (1977), 206--222. A.~M.~Vershik and S.~V.~Kerov, ``Asymptotics of the Plancherel measure of the symmetric group
and the limiting form of Young tables'', \emph{Soviet Math.\ Dokl.}\ 18 (1977), 527--531. J.~Baik, P.~Deift and K.~Johansson, ``On the distribution of the length of the longest increasing
subsequence of random permutations'', \emph{J.\ Amer.\ Math.\ Soc.}\ 12 (1999), 1119--1178. D.~Romik, \emph{The Surprising Mathematics of Longest Increasing Subsequences}, Cambridge
University Press, 2015, tells the whole story, Ulam's question included.
\item S.~Boucheron, G.~Lugosi and P.~Massart, \emph{Concentration Inequalities: A Nonasymptotic Theory of Independence}, Oxford University Press, 2013. The modern reference, with
Efron--Stein, entropy methods and many applications.
\item M.~Ledoux, \emph{The Concentration of Measure Phenomenon}, AMS, 2001. The geometric and semigroup point of view, Herbst's argument included.
\item N.~Alon and J.~H.~Spencer, \emph{The Probabilistic Method}, 4th ed., Wiley, 2016, Chapter~7. Martingales and Talagrand's inequality for combinatorics; our proof of
Theorem~\ref{thm:talagrand} follows theirs.
\item R.~Vershynin, \emph{High-Dimensional Probability}, Cambridge University Press, 2018. Concentration for data science, Johnson--Lindenstrauss included.
\end{itemize}
